% GENERATED FILE. Edit canonical lessons, then run npm run generate:books.
% canonical-source-hash: fnv1a64:e3e14416c93f04fd
% canonical-lessons: MW-C45-pet, MW-C45-angli, MW-C45-hoth, MW-C45-mudo, MW-C45-godo

\chapter{Stomach, Finger, Lip, Face, Knee}
\label{ch:mw-stomach-finger-lip-face-knee}
\hlchaptermodality{45}

\begin{chapteropening}
\textbf{By the end of this chapter:} \emph{I can say the stomach, a finger, a lip, the face, the mouth and a knee.}

Complete the last lesson of chapter 45: I can say the stomach, a finger, a lip, the face, the mouth and a knee.
\end{chapteropening}

\section[\texorpdfstring{\mw{पेट} (peṭ)}{पेट (peṭ)}]{\mw{पेट} (peṭ) — the stomach}

\label{lesson:MW-C45-pet}



\begin{quote}
Before the new one: say the Marwadi for a tooth, then the Marwadi for the throat.
\end{quote}

\subsection*{You'll want to know: \mw{पेट}}
\textbf{\mw{पेट}} — \emph{peṭ} — \textquotedblleft{}the stomach\textquotedblright{}.

\begin{cousinweb}
\textbf{\mw{पेट}} (\emph{peṭ}) is a tadbhava, a word inherited and reshaped by sound change. It comes from \emph{*peṭṭa}, \textquotedblleft{}belly\textquotedblright{}, a form CDIAL reconstructs from its descendants. CDIAL lists Old Marwari \emph{peṭa} under it (CDIAL 8376). CDIAL notes that the word may be the same as one for \textquotedblleft{}basket\textquotedblright{}.
\end{cousinweb}

\begin{grammarlens}[title={What you've built}]
One more word for this chapter.
\end{grammarlens}

\subsection*{Your turn}
\begin{itemize}
\raggedright
  \item \emph{Say it:} \emph{peṭ}
  \item \emph{Say it:} \emph{peṭ}, once more
  \item \emph{Say it:} say \emph{peṭ}
  \item \emph{Recall:} say the Marwadi for a tooth, then the Marwadi for the throat, then say \emph{peṭ} again
\end{itemize}

\begin{tcolorbox}[breakable,colback=teal!4,colframe=teal!35!black,title={Before you move on}]
What does \mw{पेट} mean? (\textquotedblleft{}The stomach\textquotedblright{}.) Say it once more.
\end{tcolorbox}
\section[\texorpdfstring{\mw{आंगळी} (āṅgḷī)}{आंगळी (āṅgḷī)}]{\mw{आंगळी} (āṅgḷī) — a finger}

\label{lesson:MW-C45-angli}



\begin{quote}
Before the new one: say the Marwadi for the throat, then the Marwadi for the stomach.
\end{quote}

\subsection*{You'll want to know: \mw{आंगळी}}
\textbf{\mw{आंगळी}} — \emph{āṅgḷī} — \textquotedblleft{}a finger\textquotedblright{}.

\begin{cousinweb}
\textbf{\mw{आंगळी}} (\emph{āṅgḷī}) is a tadbhava, a word inherited from Sanskrit and reshaped by sound change. It comes from \emph{aṅguli}, \textquotedblleft{}finger, toe\textquotedblright{}. CDIAL gives no Marwari form in this entry (CDIAL 135); the match is through Gujarati \emph{ā̃gḷī}. Marwadi, like Gujarati, says the \emph{l} curled back.
\end{cousinweb}

\begin{grammarlens}[title={What you've built}]
One more word for this chapter.
\end{grammarlens}

\subsection*{Your turn}
\begin{itemize}
\raggedright
  \item \emph{Say it:} \emph{āṅgḷī}
  \item \emph{Say it:} \emph{āṅgḷī}, once more
  \item \emph{Say it:} say \emph{āṅgḷī}
  \item \emph{Recall:} say the Marwadi for the throat, then the Marwadi for the stomach, then say \emph{āṅgḷī} again
\end{itemize}

\begin{tcolorbox}[breakable,colback=teal!4,colframe=teal!35!black,title={Before you move on}]
What does \mw{आंगळी} mean? (\textquotedblleft{}A finger\textquotedblright{}.) Say it once more.
\end{tcolorbox}
\section[\texorpdfstring{\mw{होठ} (hoṭh)}{होठ (hoṭh)}]{\mw{होठ} (hoṭh) — a lip}

\label{lesson:MW-C45-hoth}



\begin{quote}
Before the new one: say the Marwadi for the stomach, then the Marwadi for a finger.
\end{quote}

\subsection*{You'll want to know: \mw{होठ}}
\textbf{\mw{होठ}} — \emph{hoṭh} — \textquotedblleft{}a lip\textquotedblright{}.

\begin{cousinweb}
\textbf{\mw{होठ}} (\emph{hoṭh}) is a tadbhava of Sanskrit \emph{oṣṭha}, \textquotedblleft{}lip\textquotedblright{}, under which CDIAL lists Hindi and Gujarati \emph{hoṭh} (CDIAL 2563). Where the \emph{h-} came from is not settled: a proposal CDIAL records traces it to \emph{*adhaoṣṭha}, \textquotedblleft{}lower lip\textquotedblright{} (CDIAL 245a).
\end{cousinweb}

\begin{grammarlens}[title={What you've built}]
One more word for this chapter.
\end{grammarlens}

\subsection*{Your turn}
\begin{itemize}
\raggedright
  \item \emph{Say it:} \emph{hoṭh}
  \item \emph{Say it:} \emph{hoṭh}, once more
  \item \emph{Say it:} say \emph{hoṭh}
  \item \emph{Recall:} say the Marwadi for the stomach, then the Marwadi for a finger, then say \emph{hoṭh} again
\end{itemize}

\begin{tcolorbox}[breakable,colback=teal!4,colframe=teal!35!black,title={Before you move on}]
What does \mw{होठ} mean? (\textquotedblleft{}A lip\textquotedblright{}.) Say it once more.
\end{tcolorbox}
\section[\texorpdfstring{\mw{मूंडो} (mū̃ḍo)}{मूंडो (mū̃ḍo)}]{\mw{मूंडो} (mū̃ḍo) — the face, the mouth}

\label{lesson:MW-C45-mudo}



\begin{quote}
Before the new one: say the Marwadi for a finger, then the Marwadi for a lip.
\end{quote}

\subsection*{You'll want to know: \mw{मूंडो}}
\textbf{\mw{मूंडो}} — \emph{mū̃ḍo} — \textquotedblleft{}the face, the mouth\textquotedblright{}.

\begin{cousinweb}
CDIAL lists Old Marwari \emph{mūḍho}, \textquotedblleft{}face\textquotedblright{}, under Sanskrit \emph{mūrdhan}, \textquotedblleft{}head, top\textquotedblright{} (CDIAL 10247). \textbf{\mw{मूंडो}} (\emph{mū̃ḍo}) is a close form of the same word, with a nasal vowel and no breath after the \emph{ḍ}; CDIAL does not list that exact shape.
\end{cousinweb}

\begin{grammarlens}[title={What you've built}]
One more word for this chapter.
\end{grammarlens}

\subsection*{Your turn}
\begin{itemize}
\raggedright
  \item \emph{Say it:} \emph{mū̃ḍo}
  \item \emph{Say it:} \emph{mū̃ḍo}, once more
  \item \emph{Say it:} say \emph{mū̃ḍo}
  \item \emph{Recall:} say the Marwadi for a finger, then the Marwadi for a lip, then say \emph{mū̃ḍo} again
\end{itemize}

\begin{tcolorbox}[breakable,colback=teal!4,colframe=teal!35!black,title={Before you move on}]
What does \mw{मूंडो} mean? (\textquotedblleft{}The face, the mouth\textquotedblright{}.) Say it once more.
\end{tcolorbox}
\section[\texorpdfstring{\mw{गोडो} (goḍo)}{गोडो (goḍo)}]{\mw{गोडो} (goḍo) — a knee}

\label{lesson:MW-C45-godo}



\begin{quote}
Before the new one: say the Marwadi for a lip, then the Marwadi for the face.
\end{quote}

\subsection*{You'll want to know: \mw{गोडो}}
\textbf{\mw{गोडो}} — \emph{goḍo} — \textquotedblleft{}a knee\textquotedblright{}.

\begin{cousinweb}
\textbf{\mw{गोडो}} (\emph{goḍo}) is a tadbhava, a word inherited and reshaped by sound change. It comes from \emph{*goḍḍa}, \textquotedblleft{}foot, leg, knee\textquotedblright{}, a form CDIAL reconstructs from its descendants. CDIAL gives no Marwari form in this entry (CDIAL 4272); the match is through Punjabi \emph{goḍḍā} and Hindi \emph{goṛ}. CDIAL sets it in a group of words for \textquotedblleft{}heel, ankle, knee, wrist\textquotedblright{}.
\end{cousinweb}

\begin{grammarlens}[title={What you've built}]
One more word for this chapter.
\end{grammarlens}

\subsection*{Your turn}
\begin{itemize}
\raggedright
  \item \emph{Say it:} \emph{goḍo}
  \item \emph{Say it:} \emph{goḍo}, once more
  \item \emph{Say it:} say \emph{goḍo}
  \item \emph{Recall:} say the Marwadi for a lip, then the Marwadi for the face, then say \emph{goḍo} again
\end{itemize}

\begin{tcolorbox}[breakable,colback=teal!4,colframe=teal!35!black,title={Before you move on}]
What does \mw{गोडो} mean? (\textquotedblleft{}A knee\textquotedblright{}.) Say it once more.
\end{tcolorbox}
